% !TeX root = ../复分析笔记.tex
% 由主文档通过 \input 载入；请编译主文档。
\setcounter{chapter}{9} % 保留原书章号。
\chapter{辐角原理与零点计数}
\label{chap:ca-counting}

\section{对数导数如何记录零点与极点}

求一个方程的全部根可能很困难，但判断某个区域中有多少根，往往只需研究边界。连接局部根与边界的函数是 $f'/f$，称为对数导数（logarithmic derivative）。即使 $f$ 没有全局对数，这个商在 $f$ 非零且有限处仍是单值函数。

若 $a$ 是 $f$ 的 $m$ 阶零点，局部分解 $f=(z-a)^mg$ 给出
\[
 \frac{f'}f=\frac m{z-a}+\frac{g'}g,
 \qquad g(a)\ne0.
\]
因此对数导数的留数为 $m$。若 $a$ 是 $m$ 阶极点，同样计算得到留数 $-m$。零点与极点在闭路积分中分别贡献正、负整数。

\begin{theorem}[辐角原理]
\label{thm:ca-argument}
设 $f$ 在正向分段光滑简单闭曲线 $\Gamma$ 及其内部的开邻域亚纯，且 $\Gamma$ 上没有零点和极点。记内部零点总数为 $N$、极点总数为 $P$，均按重数计算，则
\begin{equation}
 \frac1{2\pi i}\oint_\Gamma\frac{f'(z)}{f(z)}\,\mathrm dz=N-P.
 \label{eq:ca-argument}
\end{equation}
\end{theorem}
\begin{proof}
曲线及其内部构成紧集，而零点、极点在相关邻域内均孤立且不在边界，所以内部只有有限个。对 $f'/f$ 应用留数定理，每个零点贡献其重数，每个极点贡献重数的负值，即得公式。
\end{proof}

沿 $\Gamma$ 连续追踪 $f(z)$ 的辐角，式~\eqref{eq:ca-argument} 也等于像曲线 $f\circ\Gamma$ 关于原点的绕数，所以
\[
 \Delta_\Gamma\arg f=2\pi(N-P).
\]
这里是沿路径连续追踪的总辐角变化，不是终点辐角主值减去起点主值；后者在闭曲线上可能总写成零，丢失绕行信息。

\section{鲁歇定理：小扰动不改变零点总数}

\begin{theorem}[鲁歇定理]
\label{thm:ca-rouche}
设 $f,g$ 在正向分段光滑简单闭曲线 $\Gamma$ 及其内部的开邻域全纯，且边界上满足
\[
 |g(z)|<|f(z)|.
\]
则 $f$ 与 $f+g$ 在 $\Gamma$ 内的零点总数相同，均按重数计算。
\end{theorem}
\begin{proof}
令 $F_t=f+tg$，$0\le t\le1$。边界上 $|F_t|\ge|f|-|g|>0$，并且紧性给出对所有 $t,z\in[0,1]\times\Gamma$ 统一的正下界。因此
\[
 N(t)=\frac1{2\pi i}\oint_\Gamma\frac{F_t'(z)}{F_t(z)}\,\mathrm dz
\]
随 $t$ 连续；由辐角原理，它又只能取整数。连续整数值函数在区间上恒定，故 $N(0)=N(1)$。
\end{proof}

几何上，$f+tg$ 的边界像在变形时始终不穿过原点，所以绕原点的圈数不变。严格不等式应在整个边界上成立，只在部分点成立不够。定理保留的是零点总数而不是每个零点的位置，更不是每个零点都保持原来的重数。

\begin{example}[不用求根公式判断根的位置]
求 $P(z)=z^5+3z+1$ 在 $|z|<1$ 和 $1<|z|<2$ 内的零点数。
\end{example}
\begin{solve}
在 $|z|=1$ 上取主项 $f=3z$，其余项 $g=z^5+1$ 满足 $|g|\le2<3=|f|$，故单位圆内有一个零点。在 $|z|=2$ 上取主项 $z^5$，有 $|3z+1|\le7<32=|z^5|$，故半径 $2$ 的圆内有五个零点。两个边界均无零点，所以圆环内有四个零点。主项的选择随半径改变，不能机械地总选次数最高项。
\end{solve}

\begin{example}[简单根在小扰动下的稳定性]
设 $a$ 是全纯函数 $f$ 的简单零点。取足够小的圆 $|z-a|=r$，使闭圆盘内没有其他零点，并令 $m_r=\min_{|z-a|=r}|f(z)|>0$。若 $h$ 在闭圆盘附近全纯且边界上 $|h|<m_r$，则 $f+h$ 在圆内按重数恰有一个零点，因此该零点也简单。这个结论可以证明近似方程仍有唯一的邻近根，但并不直接给出数值算法的误差或收敛速度。
\end{example}

\section{开映射与局部反函数}

\begin{theorem}[开映射定理]
\label{thm:ca-open-map}
区域上的非常数全纯函数把开集映成开集。
\end{theorem}
\begin{proof}
任取 $a\in D$。由恒等定理，$f-f(a)$ 不在 $a$ 附近恒零，所以 $a$ 为其某个有限阶零点。选小圆使圆内无其他零点，令 $\delta=\min_{|z-a|=r}|f(z)-f(a)|>0$。对任何满足 $|w-f(a)|<\delta$ 的 $w$，鲁歇定理说明 $f(z)-w$ 在圆内至少有一个零点。故 $B(f(a),\delta)$ 包含于该小圆盘的像，说明每个像点都是内点。将小圆选在任意给定开子集内，同样证明该开子集的像为开集。
\end{proof}

\begin{theorem}[全纯反函数定理]
\label{thm:ca-inverse}
若 $f$ 在 $a$ 附近全纯且 $f'(a)\ne0$，则存在 $a$ 的邻域 $U$ 和 $f(a)$ 的邻域 $V$，使 $f:U\to V$ 一一对应，反函数全纯，且
\[
 (f^{-1})'(w)=\frac1{f'(f^{-1}(w))}.
\]
\end{theorem}
\begin{proof}
全纯函数作为实映射光滑，且雅可比行列式为 $|f'(a)|^2>0$。由实反函数定理，可在足够小的邻域得到连续可微反函数。固定 $w_0=f(z_0)$，令 $w=f(z)$，由反函数连续性，$w\to w_0$ 时 $z\to z_0$，于是
\[
 \frac{f^{-1}(w)-f^{-1}(w_0)}{w-w_0}
 =\frac{z-z_0}{f(z)-f(z_0)}\longrightarrow\frac1{f'(z_0)}.
\]
缩小邻域可保证导数不为零，故反函数在邻域处处复可导。
\end{proof}

反过来，单射全纯函数的导数处处不为零。若在某点 $a$ 有 $f'(a)=0$，则 $f-f(a)$ 的零点阶 $m\ge2$。取小圆使其内部除 $a$ 外既无该函数的零点，也无 $f'$ 的零点；对足够接近但不同于 $f(a)$ 的 $w$，鲁歇定理给出 $f-w$ 在圆内有 $m$ 个零点。它们不等于 $a$，导数又不为零，所以均简单，因而至少有两个不同原像，违背单射。这一步解释了共形映射为什么需要排除临界点。

\section{无零点函数何时有全纯对数}

\begin{theorem}[单连通区域上的全纯对数]
\label{thm:ca-holomorphic-log}
若 $D$ 单连通，$f$ 在 $D$ 全纯且处处非零，则存在全纯函数 $L$ 使 $e^L=f$。任意两个这样的 $L$ 相差 $2k\pi i$，其中 $k$ 是固定整数。因而 $f$ 在 $D$ 上也存在任意正整数阶的全纯根。
\end{theorem}
\begin{proof}
$f'/f$ 在 $D$ 全纯，由单连通形式的柯西定理存在原函数 $G$。计算
\[
 (fe^{-G})'=e^{-G}(f'-fG')=0,
\]
故 $fe^{-G}=c\ne0$。任选满足 $e^\ell=c$ 的一个复数 $\ell$，取 $L=G+\ell$ 即可。两个对数之差连续且只取 $2\pi i\mathbb Z$ 中的值，由连通性为常数。最后 $e^{L/n}$ 是 $f$ 的一个全纯 $n$ 次根。
\end{proof}

一般区域上，准确的对数存在条件是所有闭路都满足 $\oint f'/f\,\mathrm dz=0$：必要性来自对数的导数，充分性来自原函数构造及上面的指数消去。单连通是容易验证的充分条件，却不是每一个具体函数所必需的条件；例如常函数 $f=1$ 在圆环上也有全纯对数零。

\begin{proposition}[全纯根的周期判据]
\label{prop:ca-root-periods}
设 $f$ 在区域 $D$ 全纯且无零点，$k$ 为正整数。存在全纯函数 $h$ 使 $h^k=f$，当且仅当对每条分段 $C^1$ 闭路 $\gamma\subset D$ 都有
\[
 \frac1{2\pi i}\oint_\gamma\frac{f'(z)}{f(z)}\,\mathrm dz\in k\mathbb Z.
\]
\end{proposition}
\begin{proof}
若 $f=h^k$，则 $h$ 处处非零且 $f'/f=k h'/h$。$\frac1{2\pi i}\oint h'/h$ 是像曲线 $h\circ\gamma$ 绕零点的整数绕数，因此条件必要。

反之，固定 $a\in D$，选 $c^k=f(a)$，对从 $a$ 到 $z$ 的分段光滑路径 $\sigma$ 定义
\[
 h(z)=c\exp\left(\frac1k\int_\sigma\frac{f'(\zeta)}{f(\zeta)}\,\mathrm d\zeta\right).
\]
两条路径的积分之差是某条闭路的积分，按假设属于 $2\pi i k\mathbb Z$，故指数值与路径无关。在任一点附近固定到该点的路径，再接一小段线段，可得 $h$ 全纯且 $h'/h=f'/(kf)$。所以 $(h^k/f)'=0$，由基点值及连通性得到 $h^k/f\equiv1$。
\end{proof}

例如在去心平面上，$f(z)=z$ 绕单位圆的对数导数积分为 $2\pi i$，故没有全纯平方根；$f(z)=z^2$ 的积分为 $4\pi i$，平方根 $h(z)=z$ 存在，但 $f$ 仍没有全纯对数。全纯对数要求所有周期为零，全纯 $k$ 次根只要求相应绕数能被 $k$ 整除。

\begin{proposition}[全纯映射的局部标准形]
\label{prop:ca-local-normal-form}
若 $f$ 在 $a$ 附近全纯，$f-f(a)$ 在 $a$ 有 $m\ge1$ 阶零点，则存在 $a$ 附近的双全纯局部坐标 $\psi$，满足 $\psi(a)=0$ 且
\[
 f(z)-f(a)=\psi(z)^m.
\]
\end{proposition}
\begin{proof}
由零点阶的分解，$f(z)-f(a)=(z-a)^m g(z)$，其中 $g$ 全纯且 $g(a)\ne0$。在一个充分小的圆盘内 $g$ 无零点，取全纯对数 $L$，令 $\psi(z)=(z-a)e^{L(z)/m}$。则 $\psi'(a)=e^{L(a)/m}\ne0$，全纯反函数定理保证缩小邻域后 $\psi$ 双全纯，代入即得公式。
\end{proof}

因此临界点附近的映射具有幂映射的局部结构。选取足够小的 $r$ 使 $B(0,r)$ 包含于坐标像，每个满足 $0<|w-f(a)|<r^m$ 的 $w$ 在 $\psi^{-1}(B(0,r))$ 中恰有 $m$ 个不同原像，因为非零数的 $m$ 次根互异。这同时解释了零点重数、局部多对一和导数消失之间的关系。

\section{习题}
\label{sec:ca-counting-ex}
\begin{enumerate}
 \item 求 $z^4+3z+1$ 在单位圆内的零点数。
 \item 求 $z^4+z+1$ 在 $|z|<2$ 内的零点数，并说明边界上是否可能有根。
 \item 对 $f(z)=(z-1)^3/(z+1)^2$，计算沿正向圆 $|z|=2$ 的 $\oint f'/f\,\mathrm dz$。
 \item 若 $f$ 在单连通区域 $D$ 全纯且无零点，构造一个全纯平方根，并说明所有平方根之间的关系。
 \item 证明非常数全纯函数不可能把一个区域映入一条直线或一个圆周。
\end{enumerate}

